%% it/sections/07_public_history.tex
%% ============================================================
%% 07_public_history.tex — Guida Public History scolastica
%% ============================================================

\chapter{Public History in Classe}
\markboth{Public History in Classe}{Public History in Classe}

\lettrine[lines=3]{G}{iocare} ai Vespri Siciliani in classe non è
«usare il gioco per insegnare la storia». È qualcosa di più ambizioso:
è mettere gli studenti nella posizione di chi deve fare scelte storiche
senza sapere come andrà a finire. Questo cambia il loro rapporto
con il passato — e con il presente.

Questa sezione è destinata agli insegnanti e ai GM-educatori che
vogliono usare \textit{Vespri 1282} in contesto scolastico. Si
rivolge principalmente alla \textbf{scuola secondaria di secondo
grado} (licei classici, scientifici, delle scienze umane, istituti
tecnici), ma molte delle indicazioni sono adattabili.

\nota{Una premessa onesta}{
  Il gioco di ruolo in classe funziona quando l'insegnante è
  disposto a cedere una parte del controllo sull'esperienza.
  Non sapete esattamente cosa diranno i vostri studenti, cosa
  decideranno, come reagiranno. Questo è il punto. Se volete
  controllare ogni variabile, questo strumento non fa per voi.
  Se volete sorprendere e sorprendervi, continuate a leggere.
}

\secsep

\section{Perché i Vespri Siciliani in Classe}

\subsection{Agganci Curriculari}

\textit{Vespri 1282} si collega direttamente a nuclei tematici
presenti nei programmi ministeriali italiani:

\begin{itemize}
  \item \textbf{Storia medievale} (II anno, biennio): il Mediterraneo
    come spazio di conflitto e contatto; il Papato e l'Impero;
    le monarchie europee del XIII secolo
  \item \textbf{Geopolitica medievale}: la competizione tra Angiò,
    Aragona, Bisanzio; la crociata come strumento politico
  \item \textbf{Storia della Sicilia}: dall'emirato normanno
    agli Svevi agli Angioini; la Sicilia come spazio multiculturale
  \item \textbf{Educazione civica}: potere e resistenza; diritti
    e oppressione; la violenza collettiva come risposta politica
  \item \textbf{Storiografia}: il dibattito tra «storia degli
    eventi» e «storia delle strutture»; le fonti primarie
    e i loro limiti
\end{itemize}

\subsection{Competenze Trasversali}

Oltre ai contenuti disciplinari, \textit{Vespri 1282} sviluppa:

\begin{itemize}
  \item \textbf{Pensiero critico}: valutare informazioni
    parziali e contraddittorie
  \item \textbf{Empatia storica}: assumere il punto di vista
    di persone in contesti molto diversi dal proprio
  \item \textbf{Cooperazione}: lavorare con un gruppo verso
    obiettivi condivisi, gestire conflitti interni
  \item \textbf{Comunicazione}: argomentare, persuadere,
    negoziare in tempo reale
  \item \textbf{Pensiero controfattuale}: immaginare alternative
    storiche plausibili e valutarne le conseguenze
\end{itemize}

\secsep

\section{Struttura dell'Attività}

\subsection{Tempi Consigliati}

\begin{tabular}{lll}
\toprule
\cinzel Fase & \cinzel Durata & \cinzel Quando \\
\midrule
Preparazione storica & 2--3 ore lezione & settimane prima \\
Presentazione regole & 1 ora lezione & giorno prima o mattina stessa \\
One-shot & 3--4 ore & blocco pomeridiano o giornata \\
Debriefing & 1--2 ore & immediatamente dopo \\
Elaborato scritto & a casa & settimana successiva \\
\bottomrule
\end{tabular}

\subsection{Preparazione Storica}

Prima della sessione di gioco, gli studenti devono avere
una conoscenza di base del contesto. Si consiglia:

\begin{enumerate}
  \item \textbf{Lettura guidata} dell'Introduzione Storica
    di questo manuale (sezione I), in classe o a casa
  \item \textbf{Analisi di una fonte primaria}: un brano
    di Giovanni Villani o Bartolomeo di Neocastro sulle
    cause della rivolta
  \item \textbf{Discussione in classe}: \textit{Chi aveva
    torto marcio? Chi aveva ragione parziale? Chi non
    aveva scelta?}
  \item \textbf{Assegnazione dei ruoli}: scegliere o
    assegnare le fazioni prima della sessione, così
    gli studenti sanno chi «sono» e possono prepararsi
\end{enumerate}

\nota{Sull'assegnazione dei ruoli}{
  Per classi numerose (30+ studenti), dividete in quattro
  gruppi e assegnate le fazioni. Se possibile, lasciate
  che ogni gruppo scelga i propri playbook internamente.
  Il processo di scelta è già un'attività: \textit{chi
  vuoi essere in questa storia?}
}

\subsection{Il Giorno della Sessione}

\textbf{Mattina (se è una giornata intera):}
\begin{itemize}[nosep]
  \item 9:00--10:00 — Ripasso storico, presentazione
    delle regole, domande
  \item 10:00--10:15 — Pausa, i GM si posizionano ai tavoli
  \item 10:15--10:30 — Introduzione in-world: ogni GM
    legge il testo di apertura della propria fazione
\end{itemize}

\textbf{Sessione di gioco:}
\begin{itemize}[nosep]
  \item 10:30--12:00 — Prima parte (fino alla Voce)
  \item 12:00--13:00 — Pranzo
  \item 13:00--14:30 — Seconda parte (fino alla Crisi)
  \item 14:30--15:30 — Terza parte (fino alle Campane)
  \item 15:30--16:00 — Risoluzione e rivelazione dell'esito
\end{itemize}

\secsep

\section{Testi di Apertura per Fazione}
\textit{(Da leggere ad alta voce dal GM all'inizio della sessione)}

\subsection{Apertura: I Congiurati}

\begin{citbox}
\textit{Palermo, febbraio 1282. La notte è fredda. Voi siete qui —
in questo posto, in questo momento — perché qualcuno si è fidato
di voi abbastanza da dirvi la verità.}

\textit{La verità è questa: tra poche settimane, se tutto va come
deve andare, Palermo non sarà più angioina. La verità è anche
questa: non tutto va come deve andare. E la verità più scomoda
è che nessuno di voi sa esattamente cosa succederà — solo quello
che \textbf{voi} farete.}

\textit{Avete scelto di essere qui. Adesso sceglierete il resto.}
\end{citbox}

\subsection{Apertura: I Baroni}

\begin{citbox}
\textit{Palermo, febbraio 1282. Avete ricevuto un messaggio che
non avrebbero dovuto mandarvi. O forse sì — forse qualcuno
ha deciso che era ora di sapere.}

\textit{La proposta è semplice: c'è qualcosa che si muove.
Qualcosa di grosso. Potete farne parte — o potete restare
a guardare e sperare che non vi travolga. Ma il messaggio
dice anche un'altra cosa: il tempo per decidere è adesso.
Non domani. Non la settimana prossima.}

\textit{Adesso.}
\end{citbox}

\subsection{Apertura: Il Clero}

\begin{citbox}
\textit{Palermo, febbraio 1282. Le campane hanno appena suonato
Compieta. La città si è fatta quieta. Ma la quiete non è pace —
è la quiete di chi trattiene il respiro.}

\textit{Voi servite Dio. Servite Roma. Servite il vostro popolo.
Stanotte, per la prima volta, non siete sicuri che queste tre
cose vadano nella stessa direzione. Stanotte potreste fare
qualcosa che non si può disfare.}

\textit{Pregate. E poi decidete.}
\end{citbox}

\subsection{Apertura: Il Popolo}

\begin{citbox}
\textit{Palermo, febbraio 1282. È mattina. Il mercato è aperto.
I soldati francesi sono già in piazza — come sempre, da sedici
anni, come sempre.}

\textit{Voi non sapete niente di nessuna cospirazione. Sapete
solo quello che sapete: le tasse sono arrivate di nuovo. Tuo
cugino è stato picchiato la settimana scorsa senza motivo.
Il prezzo del pane è salito ancora. E qualcuno vi ha detto,
ieri sera, che tra qualche settimana potrebbe cambiare tutto.}

\textit{Non sapete se ci credete. Ma vi hanno chiesto di esserci.
E voi siete qui.}
\end{citbox}

\secsep

\section{Debriefing Strutturato}

Il debriefing è la parte più importante dell'attività didattica.
Senza debriefing, il gioco è solo un'esperienza. Con il debriefing,
diventa un'occasione di apprendimento profondo.

\subsection{Prima Fase: De-roll (10--15 minuti)}

Prima di discutere storia o significati, è fondamentale aiutare
gli studenti a uscire dai personaggi.

\textbf{Domande di de-roll:}
\begin{itemize}[nosep]
  \item «Cosa stava provando il tuo personaggio nell'ultima scena?»
  \item «Cosa stavi provando \textbf{tu} in quella stessa scena?»
  \item «C'è stata una scelta che hai fatto e di cui sei
    fiero/a? Una di cui ti sei pentito/a?»
\end{itemize}

Questo passaggio è importante soprattutto se il gioco ha prodotto
momenti emotivamente intensi. Non saltarlo.

\subsection{Seconda Fase: Riflessione Storica (30--45 minuti)}

\textbf{Domande per la discussione:}

\begin{enumerate}
  \item \textit{Il vostro esito corrisponde a quello storico?
    Se no, cosa è andato diversamente? Perché?}

  \item \textit{Quali personaggi storici avete incontrato?
    Come li avete interpretati? Corrispondevano a quello
    che sapevate di loro?}

  \item \textit{C'è stato un momento in cui avete capito
    qualcosa della situazione storica che prima non
    avevate capito? Cosa?}

  \item \textit{La rivolta dei Vespri era «giusta»?
    Era «inevitabile»? Queste sono due domande diverse —
    discutiamole separatamente.}

  \item \textit{Cosa avrebbe potuto fare Carlo d'Angiò
    per evitare la rivolta? Era possibile?}
\end{enumerate}

\subsection{Terza Fase: Connessioni al Presente (15--20 minuti)}

\begin{enumerate}
  \item \textit{Ci sono situazioni nel mondo contemporaneo
    che ricordano la struttura di potere che avete giocato?
    (Attenzione: non cercate analogie dirette — cercate
    strutture simili)}

  \item \textit{Chi scrive la storia dei Vespri? Le fonti
    che avete letto — Villani, Bartolomeo di Neocastro —
    chi erano? Perché scrivevano?}

  \item \textit{Il Popolo ha vinto. Eppure i loro nomi
    non sono nelle cronache. Com'è possibile? È giusto?}
\end{enumerate}

\secsep

\section{Elaborato Scritto}

\subsection{Tipologie di Elaborato (a scelta dello studente)}

\textbf{A — Il diario del personaggio}\\
Scrivi il diario del tuo personaggio per le tre settimane
precedenti il 30 marzo 1282. Include: le sue motivazioni,
le sue paure, le scelte che ha fatto durante la sessione,
e cosa pensa che accadrà dopo.

\textbf{Requisiti:} 600--800 parole. Coerenza storica (nessun
anacronismo). Almeno un riferimento a una fonte storica
reale (citata in bibliografia).

\filetto

\textbf{B — L'analisi storiografica}\\
Confronta due interpretazioni della «cospirazione aragonese»:
quella di Runciman (complotto organizzato) e quella di Abulafia
(rivolta spontanea). Quale ti sembra più convincente, sulla
base di quello che hai studiato e di quello che hai vissuto
in gioco? Argomenta.

\textbf{Requisiti:} 800--1000 parole. Struttura argomentativa
(tesi, argomenti, contro-argomenti, conclusione). Almeno due
fonti bibliografiche.

\filetto

\textbf{C — La fonte alternativa}\\
Immagina di essere uno storico del XIV secolo che scrive
la propria cronaca dei Vespri Siciliani. Il tuo punto di
vista è diverso da quello di Villani: sei un frate siciliano,
un mercante genovese, un funzionario angioino scampato alla
notte del 30 marzo. Scrivi la tua cronaca.

\textbf{Requisiti:} 700--900 parole. Stile coerente con la
scrittura medievale (prima persona, riferimenti religiosi,
linguaggio formale ma non accademico). Nota finale dell'autore
che spiega le scelte fatte.

\secsep

\section{Note per l'Insegnante}

\subsection{Gestione delle Emozioni}

Il gioco può produrre momenti emotivamente intensi — soprattutto
al tavolo del Popolo, dove i personaggi sono vulnerabili e le
conseguenze sono concrete. Alcune indicazioni:

\begin{itemize}
  \item Stabilite una \textbf{parola di sicurezza} prima della
    sessione: una parola che chiunque può dire per fermare
    il gioco senza spiegazioni. Rispettatela senza eccezioni.
  \item Il GM di ogni tavolo deve monitorare lo stato emotivo
    dei giocatori. Se qualcuno è chiaramente a disagio,
    offritegli la possibilità di fare una pausa.
  \item La violenza nel gioco deve restare \textbf{narrativa}
    — mai reenactment fisico. I combattimenti si risolvono
    con i dadi, non con i corpi.
\end{itemize}

\subsection{Adattamenti per Classi Diverse}

\textbf{Classi con studenti con BES/DSA:}
\begin{itemize}[nosep]
  \item Fornite le schede personaggio con testo semplificato
  \item Permettete la consultazione del manuale durante il gioco
  \item Assegnate i ruoli con più piattaforme dialogiche
    (es. Notaio, Canonico) agli studenti che preferiscono
    non essere al centro dell'attenzione
\end{itemize}

\textbf{Classi numerose (35+ studenti):}
\begin{itemize}[nosep]
  \item Aggiungete un quinto «tavolo osservatori» che seguono
    una fazione e prendono appunti, poi riferiscono durante
    il debriefing
  \item Riducete ogni fazione a 5--6 personaggi, eliminando
    i playbook meno centrali alla narrazione
\end{itemize}

\textbf{Tempo ridotto (2 ore invece di 4):}
\begin{itemize}[nosep]
  \item Eliminate il terzo atto (Le Campane) e rivelate
    l'esito basandovi solo sui punteggi dopo la Crisi
  \item Focalizzate ogni fazione su un singolo obiettivo
    principale invece di tutti
\end{itemize}

\subsection{Raccordo con le Indicazioni Nazionali}

\textit{Vespri 1282} si collega ai seguenti obiettivi delle
Indicazioni Nazionali per i Licei (DM 211/2010):

\begin{itemize}[nosep]
  \item Comprendere il cambiamento e la diversità dei tempi
    storici in una dimensione diacronica e sincronica
  \item Collocare l'esperienza personale in un sistema di
    regole fondato sul reciproco riconoscimento dei diritti
  \item Orientarsi nel tessuto produttivo del proprio territorio
    (per le classi dei tecnici: comprensione delle strutture
    economiche del medioevo meridionale)
\end{itemize}

\filettodoppio